\documentclass[10pt,a4paper]{amsart}
\usepackage{amsmath,amssymb,amsthm,graphicx}
% Retained preamble for the Stacks Project chapter snapshot in this directory.
% The chapter's own preamble.tex is a copy of the shared Stacks Project preamble; the block
% below is sliced MECHANICALLY, line for line, from it. Nothing was added, removed or
% redefined; the wrapper class, the xr-hyper/hyperref wiring and the externaldocument list
% are omitted because the assembler supplies the document class.
%
% lines taken from preamble.tex: 3, 5, 9, 11, 12, 156, 157, 158, ... (51 lines total)
%
% package coverage restored from the chapter's own preamble (2026-10-13): xy, xy-2cell

\usepackage[all]{xy}
\xyoption{2cell}
\UseAllTwocells

\newenvironment{reference}{\comment}{\endcomment}
\newenvironment{slogan}{\comment}{\endcomment}
\newenvironment{history}{\comment}{\endcomment}
\theoremstyle{plain}
\newtheorem{theorem}[subsection]{Theorem}
\newtheorem{proposition}[subsection]{Proposition}
\newtheorem{lemma}[subsection]{Lemma}
\theoremstyle{definition}
\newtheorem{definition}[subsection]{Definition}
\newtheorem{example}[subsection]{Example}
\newtheorem{exercise}[subsection]{Exercise}
\newtheorem{situation}[subsection]{Situation}
\theoremstyle{remark}
\newtheorem{remark}[subsection]{Remark}
\newtheorem{remarks}[subsection]{Remarks}
\numberwithin{equation}{subsection}
\def\lim{\mathop{\mathrm{lim}}\nolimits}
\def\colim{\mathop{\mathrm{colim}}\nolimits}
\def\Spec{\mathop{\mathrm{Spec}}}
\def\Hom{\mathop{\mathrm{Hom}}\nolimits}
\def\Ext{\mathop{\mathrm{Ext}}\nolimits}
\def\SheafHom{\mathop{\mathcal{H}\!\mathit{om}}\nolimits}
\def\SheafExt{\mathop{\mathcal{E}\!\mathit{xt}}\nolimits}
\def\Sch{\mathit{Sch}}
\def\Mor{\mathop{\mathrm{Mor}}\nolimits}
\def\Ob{\mathop{\mathrm{Ob}}\nolimits}
\def\Sh{\mathop{\mathit{Sh}}\nolimits}
\def\NL{\mathop{N\!L}\nolimits}
\def\CH{\mathop{\mathrm{CH}}\nolimits}
\def\proetale{{pro\text{-}\acute{e}tale}}
\def\etale{{\acute{e}tale}}
\def\QCoh{\mathit{QCoh}}
\def\Ker{\mathop{\mathrm{Ker}}}
\def\Im{\mathop{\mathrm{Im}}}
\def\Coker{\mathop{\mathrm{Coker}}}
\def\Coim{\mathop{\mathrm{Coim}}}
\DeclareMathSymbol{\boxtimes}{\mathbin}{AMSa}{"02}
\def\QCohstack{\mathcal{QC}\!\mathit{oh}}
\def\Cohstack{\mathcal{C}\!\mathit{oh}}
\def\Spacesstack{\mathcal{S}\!\mathit{paces}}
\def\Quotfunctor{\mathrm{Quot}}
\def\Hilbfunctor{\mathrm{Hilb}}
\def\Curvesstack{\mathcal{C}\!\mathit{urves}}
\def\Polarizedstack{\mathcal{P}\!\mathit{olarized}}
\def\Complexesstack{\mathcal{C}\!\mathit{omplexes}}
\def\Pic{\mathop{\mathrm{Pic}}\nolimits}
\def\Picardstack{\mathcal{P}\!\mathit{ic}}
\def\Picardfunctor{\mathrm{Pic}}
\def\Deformationcategory{\mathcal{D}\!\mathit{ef}}
\usepackage{cleveref}
\crefname{originallemma}{original lemma}{original lemmas}
\Crefname{originallemma}{Original lemma}{Original lemmas}
\makeatletter
\newlabel{lemma-stalk-of-module-with-support}{{0AET}{source}{}{sourcepointer-1}{}}
\newlabel{lemma-stalk-of-module-with-support@cref}{{[originallemma][0][]0AET}{[1][0][]source}{}{}{}}
\makeatother
\crefname{originallemma}{original lemma}{original lemmas}
\Crefname{originallemma}{Original lemma}{Original lemmas}
\makeatletter
\newlabel{lemma-stalk-pushforward-with-support}{{0AEQ}{source}{}{sourcepointer-2}{}}
\newlabel{lemma-stalk-pushforward-with-support@cref}{{[originallemma][0][]0AEQ}{[1][0][]source}{}{}{}}
\makeatother
\crefname{originallemma}{original lemma}{original lemmas}
\Crefname{originallemma}{Original lemma}{Original lemmas}
\makeatletter
\newlabel{more-algebra-lemma-neighbourhood-isomorphism}{{05E9}{source}{}{sourcepointer-3}{}}
\newlabel{more-algebra-lemma-neighbourhood-isomorphism@cref}{{[originallemma][0][]05E9}{[1][0][]source}{}{}{}}
\makeatother
\title{Lf\textasciicircum{}*Rf\_*Q = Q for Q supported on f\textasciicircum{}\{-1\}Z when f is an isomorphism over Z and flat along it}
\author{The Stacks Project Authors}
\date{stacks-project:spaces-pushouts.tex:local-snapshot-20261012; GNU FDL 1.2 or later}
\begin{document}
\maketitle
\noindent\textit{Editorial note.} Statement span is lines 1993-2000 of sources/stacks-spaces-pushouts/pinned-original-source.tex; proof span is lines 2002-2043, the complete author proof (40 body lines).\par\medskip
\noindent\textit{Editorial note.} Navigation stubs: 3 cited-but-not-included labels are supplied as non-load-bearing pointers with their pinned tags where the snapshot tag index has them.\par\medskip
\noindent\textit{Editorial note.} extra\_packages is empty on purpose: the retained preamble already loads xy with options and a second declaration clashes in the assembler wrapper.\par\medskip
\noindent\textit{Editorial note.} Text is retained verbatim, including the author's own typos and placeholders.\par\medskip
\noindent\textit{Editorial note.} Queue index 14 of queue\_spaces-pushouts.json (pb=1808, stmt 1993-2000, proof 2002-2043).\par\medskip
\noindent\textit{Editorial note.} spaces-pushouts wave-1: fresh adjudication accept (2026-10-14); six-check dispatched pool (§2 row appended 02:4x).\par\medskip
\noindent\textit{Allowed standard prerequisites (editorial).}\begin{itemize}
\item Explicit prerequisite boundary. The cited inputs are prerequisites; the argument in the retained span is not replaced by any cited result.
\end{itemize}\medskip
\ifdefined\hypertarget\hypertarget{sourcepointer-1}{}\fi
\noindent\textit{Source reference.} Original lemma 0AET: sources/stacks-spaces-pushouts/pinned-original-source.tex, source lines 1957--1967. Cited by the retained statement or proof but not part of the collected spans; supplied only as a navigation pointer. The argument does not depend on its content.\par\medskip
\ifdefined\hypertarget\hypertarget{sourcepointer-2}{}\fi
\noindent\textit{Source reference.} Original lemma 0AEQ: sources/stacks-spaces-pushouts/pinned-original-source.tex, source lines 1840--1850. Cited by the retained statement or proof but not part of the collected spans; supplied only as a navigation pointer. The argument does not depend on its content.\par\medskip
\ifdefined\hypertarget\hypertarget{sourcepointer-3}{}\fi
\noindent\textit{Source reference.} Original lemma 05E9: sources/stacks-spaces-pushouts/pinned-original-source.tex, source lines 2043--2043. Cited by the retained statement or proof but not part of the collected spans; supplied only as a navigation pointer. The argument does not depend on its content.\par\medskip
\begin{lemma}
\label{lemma-formal-glueing-on-closed}
Let $S$ be a scheme. Let $f : Y \to X$ be a morphism of algebraic spaces
over $S$. Let $Z \subset X$ be a closed subspace. Assume $f^{-1}Z \to Z$
is an isomorphism and that $f$ is flat in every point of $f^{-1}Z$. For any
$Q$ in $D_\QCoh(\mathcal{O}_Y)$ supported on $|f^{-1}Z|$ we have
$Lf^*Rf_*Q = Q$.
\end{lemma}

\begin{proof}
We show the canonical map $Lf^*Rf_*Q \to Q$ is an isomorphism
by checking on stalks at $\overline{y}$. If $\overline{y}$ is not
in $f^{-1}Z$, then both sides are zero and the result is true.
Assume the image $\overline{x}$ of $\overline{y}$ is in $Z$.
By Lemma \ref{lemma-stalk-pushforward-with-support} we have
$Rf_*Q_{\overline{x}} = Q_{\overline{y}}$ and since $f$ is flat
at $\overline{y}$ we see that
$$
(Lf^*Rf_*Q)_{\overline{y}} =
(Rf_*Q)_{\overline{x}}
\otimes_{\mathcal{O}_{X, \overline{x}}}
\mathcal{O}_{Y, \overline{y}} =
Q_{\overline{y}} \otimes_{\mathcal{O}_{X, \overline{x}}}
\mathcal{O}_{Y, \overline{y}}
$$
Thus we have to check that the canonical map
$$
Q_{\overline{y}} \otimes_{\mathcal{O}_{X, \overline{x}}}
\mathcal{O}_{Y, \overline{y}}
\longrightarrow Q_{\overline{y}}
$$
is an isomorphism in the derived category. Let
$I_{\overline{x}} \subset \mathcal{O}_{X, \overline{x}}$ be the
stalk of the ideal sheaf defining $Z$. Since $Z \to X$ is locally of
finite presentation this ideal is finitely generated and the
cohomology groups of $Q_{\overline{y}}$
are $I_{\overline{y}} = I_{\overline{x}}\mathcal{O}_{Y, \overline{y}}$-power
torsion by Lemma \ref{lemma-stalk-of-module-with-support} applied to $Q$ on $Y$.
It follows that they are also $I_{\overline{x}}$-power torsion.
The ring map
$\mathcal{O}_{X, \overline{x}} \to \mathcal{O}_{Y, \overline{y}}$
is flat and induces an isomorphism after dividing by
$I_{\overline{x}}$ and $I_{\overline{y}}$ because we assumed
that $f^{-1}Z \to Z$ is an isomorphism. Hence we see that
the cohomology modules of
$Q_{\overline{y}} \otimes_{\mathcal{O}_{X, \overline{x}}}
\mathcal{O}_{Y, \overline{y}}$
are equal to the cohomology modules of $Q_{\overline{y}}$ by
More on Algebra, Lemma \ref{more-algebra-lemma-neighbourhood-isomorphism}
which finishes the proof.
\end{proof}

\end{document}
